\documentclass{article}
\usepackage{amsmath, amssymb, amsthm}
\usepackage{hyperref}
\usepackage{geometry}
\geometry{margin=1in}

\title{Erd\H{o}s Problem \#1210}
\date{Last updated: 2026-04-04}
\author{Source: erdosproblems.com}

\begin{document}
\maketitle

\noindent\textbf{Status:} OPEN \quad \textbf{No prize}

\noindent\textbf{Tags:} number theory

\medskip
\noindent\textbf{Problem Statement:}

Let $A\subseteq [1,n)$ be a set of integers such that $(a,b)=1$ for all distinct $a,b\in A$. Is it true that\[\sum_{a\in A}\frac{1}{n-a}\leq \sum_{p<n}\frac{1}{p}+O(1)?\]

\bigskip
\noindent\textbf{Remarks:}

In \cite{Er80} he claims he 'did not state [this] quite correctly' in \cite{Er77c}. The problem in \cite{Er77c} which Erd\H{o}s is presumably referring to states that if $n<q_1<\cdots<q_k\leq m$ is the set of primes in $(n,m]$ then\[\sum \frac{1}{q_i-n}< \sum_{p<m-n}\frac{1}{p}+O(1).\]See also [460] and [950].

\bigskip
\noindent\textbf{References:}

[Er77c] Erd\H{o}s, Paul, Problems and results on combinatorial number theory. III. Number theory day (Proc. Conf., Rockefeller Univ.,
New York, 1976) (1977), 43-72.
 
 [Er80] Erd\H{o}s, Paul, A survey of problems in combinatorial number theory. Ann. Discrete Math. (1980), 89-115.

\bigskip
\noindent\small{Source: \url{https://www.erdosproblems.com/1210}}
\end{document}
